\section*{Chương 4: Tổng kết}
\addcontentsline{toc}{section}{Chương 4: Tổng kết}

\subsection*{1. Ưu điểm và kết quả đạt được}
Qua quá trình nghiên cứu và thực hiện đề tài, nhóm đã hoàn thành các mục tiêu đã đặt ra và rút ra một số ưu điểm như sau:
\begin{itemize}
    \item Hiểu rõ hơn bản chất lý thuyết đằng sau phép biến đổi Haar. Củng cố các kiến thức Đại số tuyến tính đã được học.
    \item Chương trình được xây dựng với giao diện rõ ràng, giúp trực quan hóa quá các quá trình nén ảnh bằng hình ảnh, thông số và đồ thị.
    \item Nhóm đã phân tích được các phương pháp và thông số đánh giá mức độ nén của một bức ảnh. Ngoài ra còn có các nhận xét sâu hơn về tính ứng dụng của phép biến đổi Haar Wavelet trong xử lý ảnh.
    \item Quá trình làm việc chung đã giúp các thành viên rèn luyện tính kỷ luật, nâng cao kỹ năng giao tiếp, giải quyết vấn đề và tinh thần trách nhiệm trong làm việc nhóm.
\end{itemize}
\subsection*{2. Nhược điểm và hướng phát triển}
Bên cạnh những kết quả tích cực, đề tài vẫn còn một số điểm giới hạn có thể tiếp tục phát triển trong tương lai:
\begin{itemize}
    \item Thuật toán và chương trình hiện tại chỉ có thể nén ảnh tĩnh và ảnh xám. Đối với ảnh động và ảnh màu cần nhiều kỹ thuật và lý thuyết hơn. 
    Tuy vậy với nền tảng sử dụng phép biến đổi Haar Wavelet, nhóm có khả năng để thực hiện các phép biến đổi phức tạp hơn cho nhiều đối tượng hơn trong tương lai.
    \item Các thông số đánh giá được sử dụng đôi khi không thể hiện chính xác chất lượng ảnh đối với ảnh có nhiều thông tin phức tạp như hiệu ứng hậu cảnh bokeh. Hướng phát triển là 
    học hỏi thêm các thông số đánh giá nâng cao để thực hiện phân tích chính xác chất lượng ảnh.
\end{itemize}